\newpage
\section{Marco Teórico}

%\subsection{¿Qué es el internet de las cosas?}
%El Internet de las Cosas (IoT) es un componente clave de la cuarta revolución industrial, que facilita la conexión entre objetos, personas y dispositivos a través de Internet. Esta tecnología interactúa y trabaja en conjunto con otras tecnologías como sensores inalámbricos, machine learning e inteligencia artificial para habilitar conceptos como ciudades inteligentes, salud inteligente, transporte inteligente, agricultura inteligente e industria avanzada, entre otros sectores que se benefician de los avances y soluciones de IoT \cite{mincomercio_2022_gua}.\hfill \break

Blockchain se describe como  \cite{aracle_2014_qu}.\hfill \break

En resumen, d.\hfill \break

\subsection{Arquitectura de los sistemas basados en Internet de las Cosas}

Para a en el diagrama \ref{diag:flujo capas general} \cite{Ávila-Camacho_Moreno-Villalba_2023}.



\subsubsection{Capa 1: Captura o Dispositivo}
 Este es eor.\hfill \break
 
 
\subsubsection{Capa 2: Almacenamiento}
La capa d.\hfill \break


\subsubsection{Capa 3: Análisis}
La.\hfill \break


\subsubsection{Capa 4: Visualización}
La .\hfill \break

\subsection{Variables importantes a medir}

\subsubsection{Clasificación de la actividad física}

El s. \hfill\break
 
\textit{
''El gasto energético que produce cada actividad física (AF) se suele medir en MET (por las siglas en inglés de metabolic equivalent of task, (equivalente metabólico de la tarea). Un MET es la cantidad de energía gastada o consumida en una actividad, definida como la proporción de múltiplos o submúltiplos de la energía gastada durante el reposo en posición sentada. Por ejemplo, una actividad que consuma el doble de energía que estando sentado y en reposo equivale a 2 MET.''} \hfill\break


\subsection{Tabla de umbrales MET para catalogar una actividad:}


    En la bús\ref{ME} es.\hfill\break

    

\subsection{Cálculo del MET a partir de acelerómetro tri-axial}
intensidad de los movimientos durante la actividad física, tal como lo señalan los autores en \cite{MET4}.\hfill\break


\begin{equation}
VM = \sqrt{(Eje X)^{2} + (Eje Y)^{2} + (Eje Z)^{2}}
\end{equation}

El rtyrt



%%%%%%%%%%%%%%%%%%%%%%%%%%%%%%%%%%%%%%%%%%%%%%%%%%%%%%%%%%%%%%%%%%%%%%%%%%%%%%%%%%%%%%%%%%%%%%%%%%%%%%%%%%%%%%%%%%%%%%%%%%%%%%%%%%%%%%%%%%%%%%%%%%%%%%%%%%

